Let $f\in\Hecke$ (\cref{I.hecke.def:test-function}), and let $(N,F)$ and the bases $B_\mu$ be as in \cref{I.spec.def:eisenstein-kernel}. Let $K_{o,f}$ be the
spectral kernels of \cref{II.kernel.def:spectral-kernel} formed with these $B_\mu$, and $\operatorname{cond}$ as in \cref{I.eis.def:conductor}. Let $S(N)$ be the set of
orbits $o$ (\cref{II.kernel.def:cuspidal-datum}) such that $\operatorname{cond}(\mu_1)\mid N$ and $\operatorname{cond}(\mu_2)\mid N$ for every $\mu=(\mu_1,\mu_2)\in o$. Then:
\begin{enumerate}[label=(\alph*)]
\item if $o$ is an orbit and some $\mu\in o$ has $\operatorname{cond}(\mu_1)\nmid N$ or $\operatorname{cond}(\mu_2)\nmid N$, then $K_{o,f}(x,y)=0$ for all
  $x,y\in G(\A)^1$;
\item the set $S(N)$ is finite.
\end{enumerate}
